\documentclass{article}
% Match the supplied logo's PDF version and allow exact poster-size math fonts.
\pdfvariable minorversion=7
\usepackage{fix-cm}
\usepackage{fontspec}
\setmainfont{Arial}
\setsansfont{Arial}
\usepackage[czech]{babel}
\usepackage[paperwidth=1189mm,paperheight=841mm,margin=22mm]{geometry}
\usepackage{graphicx,xcolor,amsmath,amssymb,booktabs,array}
\renewcommand{\familydefault}{\sfdefault}
\pagestyle{empty}
\setlength{\parindent}{0pt}
\setlength{\parskip}{5mm}
\setlength{\tabcolsep}{4mm}
\setlength{\abovedisplayskip}{4mm}
\setlength{\belowdisplayskip}{4mm}
\definecolor{kmagold}{RGB}{209,156,0}
\definecolor{kmagrey}{RGB}{198,199,196}
\definecolor{kmalight}{RGB}{238,238,238}
\definecolor{kmatext}{RGB}{34,34,34}
\definecolor{phcolor}{RGB}{120,120,120}
\newcommand{\AuthorA}{Marek Ondriáš}
\newcommand{\AuthorB}{Vitalij Atamanjuk}
\input{results.tex}
\graphicspath{{../figures/}{images/}}

% Three aligned columns guide the reader from the model to estimates and validation.
\newcommand{\topic}[2]{%
  \par\vspace{5mm}{\fontsize{29}{34}\selectfont\bfseries
  \textcolor{kmagold}{#1}\enspace #2}\par\vspace{2mm}}
\newcommand{\note}[1]{{\fontsize{18}{23}\selectfont\color{phcolor}#1\par}}
\newcommand{\pairplots}[2]{%
  \noindent\includegraphics[width=.492\linewidth]{#1}\hfill
  \includegraphics[width=.492\linewidth]{#2}\par}

\begin{document}
\color{kmatext}
\fontsize{23}{30}\selectfont

\begin{minipage}[c]{.815\textwidth}
  {\fontsize{55}{64}\selectfont\bfseries
   Vícerozměrné náhodné veličiny,\\
   generování vícerozměrných veličin\par}
  \vspace{5mm}
  {\fontsize{27}{34}\selectfont\AuthorA\enspace\textcolor{kmagrey}{|}\enspace\AuthorB}
\end{minipage}\hfill
\begin{minipage}[c]{.16\textwidth}
  \includegraphics[width=\linewidth]{kma-cmyk-cz.pdf}
\end{minipage}
\par\vspace{5mm}
{\color{kmagold}\rule{\textwidth}{2.5mm}}
\par\vspace{3mm}

\noindent
\begin{minipage}[t]{.318\textwidth}
\vspace{0pt}
\topic{01}{Model a cíl programu}
Odvodit charakteristiky spojitého vektoru $(X,Y)$,
vygenerovat závislé dvojice a porovnat výběrové odhady s teorií.
\[
 f(x,y)=\begin{cases}cxy,&1<x<y<2,\\0,&\text{jinak.}\end{cases}
 \qquad
 1=c\int_1^2\!\int_x^2 xy\,\mathrm{d}y\,\mathrm{d}x
   =c\,\frac98
 \quad\Rightarrow\quad \boxed{c=\frac89}.
\]
\pairplots{joint_density.pdf}{accepted_pairs.pdf}
\note{Vlevo sdružená hustota; vpravo počty přijatých dvojic v šestiúhelnících.
Trojúhelníková podpora vynucuje $X<Y$.}

\topic{02}{Marginální rozdělení}
Pro $1<t<2$ dostáváme integrací sdružené hustoty:
\[
\begin{aligned}
 f_X(t)&=\frac49t(4-t^2),&
 F_X(t)&=\frac{-t^4+8t^2-7}{9},\\[2mm]
 f_Y(t)&=\frac49t(t^2-1),&
 F_Y(t)&=\frac{(t^2-1)^2}{9}.
\end{aligned}
\]
Mimo $(1,2)$ jsou hustoty nulové. Distribuční funkce mají
hodnotu $0$ pro $t\leq1$ a $1$ pro $t\geq2$.
\pairplots{marginal_pdf.pdf}{marginal_cdf.pdf}

\topic{03}{Teoretické charakteristiky}
\[
 \boldsymbol\mu=\begin{pmatrix}188/135\\232/135\end{pmatrix},
 \qquad
 \boldsymbol\Sigma=\frac1{18225}
 \begin{pmatrix}1106&484\\484&851\end{pmatrix}.
\]
\textbf{Závislost:} $\operatorname{Cov}(X,Y)=\TheoryCovariance>0$.
Vyšší hodnoty $X$ se pojí s vyššími hodnotami $Y$;
nenulová kovariance vylučuje nezávislost.
\end{minipage}\hfill
\begin{minipage}[t]{.318\textwidth}
\vspace{0pt}
\topic{04}{Generování a bodové odhady}
\textbf{Zamítací metoda:} nezávisle navrhujeme $U,V\sim U(1,2)$
a $Z\sim U(0,32/9)$. Dvojici přijmeme, právě když
$U<V$ a $Z\leq\frac89UV$.

Z $\ProposalCount$ návrhů vzniklo \textbf{$n=\SampleCount$ dvojic}.
Úspěšnost $\AcceptanceRate\,\%$ odpovídá teorii $9/32=28{,}125\,\%$.
Použit je pevný seed $42$.

\topic{05}{Střední hodnoty a intervalové odhady}
\includegraphics[width=\linewidth]{mean_estimates.pdf}
\note{Body: výběrové průměry; úsečky: 95\% intervaly se Studentovým
kvantilem; přerušované čáry: teorie. Každý řádek má vlastní osu.
Pro nenormální marginály jde o asymptotické přiblížení.}

\topic{06}{Rozptyly a kovariance}
\includegraphics[width=\linewidth]{covariance_estimates.pdf}
\note{Shodná barevná škála. Na diagonále jsou rozptyly,
mimo diagonálu kovariance. Výběrová matice používá jmenovatel $n-1$.}
\[
 S=\frac1{n-1}\sum_{i=1}^n
 (\mathbf z_i-\bar{\mathbf z})(\mathbf z_i-\bar{\mathbf z})^{\mathsf T},
 \qquad \mathbf z_i=(x_i,y_i)^{\mathsf T}.
\]
{\fontsize{21}{27}\selectfont
\renewcommand{\arraystretch}{1.45}
\setlength{\tabcolsep}{3mm}
\par\noindent
\begin{tabular*}{\linewidth}{@{\extracolsep{\fill}}lrrl@{}}
\toprule
Charakteristika & Teorie & Odhad & 95\% interval\\
\midrule
$\mathrm E[X]$ & $\TheoryMeanX$ & $\SampleMeanX$ & $(\LowerMeanX;\UpperMeanX)$\\
$\mathrm E[Y]$ & $\TheoryMeanY$ & $\SampleMeanY$ & $(\LowerMeanY;\UpperMeanY)$\\
$\operatorname{Var}(X)$ & $\TheoryVarianceX$ & $\SampleVarianceX$ & $(\LowerVarianceX;\UpperVarianceX)$\\
$\operatorname{Var}(Y)$ & $\TheoryVarianceY$ & $\SampleVarianceY$ & $(\LowerVarianceY;\UpperVarianceY)$\\
$\operatorname{Cov}(X,Y)$ & $\TheoryCovariance$ & $\SampleCovariance$ & ---\\
\bottomrule
\end{tabular*}\par}

\textbf{Technologie:} Python, NumPy (generování a odhady),
SymPy (přesné integrály), SciPy (kvantily), Matplotlib a Jupyter.
\end{minipage}\hfill
\begin{minipage}[t]{.318\textwidth}
\vspace{0pt}
\topic{07}{Výběr proti teoretickému rozdělení}
\pairplots{comparison_x_pdf.pdf}{comparison_y_pdf.pdf}
\pairplots{comparison_x_cdf.pdf}{comparison_y_cdf.pdf}
Histogramy sledují teoretické hustoty. Největší absolutní odchylka
empirické distribuční funkce od teorie je
$D_X=\CdfDeviationX$, $D_Y=\CdfDeviationY$.
\note{Jde o popis odchylky v daném výběru, nikoli o výsledek testu hypotézy.}

\topic{08}{Vliv rozsahu výběru}
\pairplots{summary_mean_error.pdf}{summary_variance_error.pdf}
\pairplots{summary_mean_ci_width.pdf}{summary_variance_ci_width.pdf}
\note{Nahoře absolutní chyby; dole šířky 95\% intervalů.
Rozsahy $n=1\,000,5\,000,10\,000,25\,000,50\,000$ jsou vnořené prefixy
jednoho doplněného výběru; nejde o nezávislá opakování.}

\topic{09}{Závěry a využití}
\textbf{Více dat zmenšuje nejistotu.} Intervaly pro průměry jsou
při $n=50\,000$ přibližně $\WidthRatioX\times$ (X) a $\WidthRatioY\times$ (Y)
užší než při $n=1\,000$. Chyba jednotlivého odhadu však nemusí klesat monotónně.

\textbf{95\% neznamená jistotu.} V původním výběru oba intervaly pro
průměry těsně míjejí teoretické hodnoty. Nominální spolehlivost se
vztahuje k opakovanému výběru, nikoli k záruce pro jediný interval.

\textbf{Omezení a využití.} Intervaly pro rozptyly používají
$\chi^2$ postup určený pro normální data; zde jsou jen orientační.
Program umožňuje kontrolu simulátoru a studium odhadů závislých veličin.
\end{minipage}
\end{document}
